\documentclass[11pt,a4paper]{article}
\usepackage[margin=2.5cm]{geometry}
\usepackage{amsmath,amssymb,amsthm}
\usepackage[hidelinks]{hyperref}
\usepackage{booktabs}
\usepackage{microtype}
\usepackage{lineno}

\newtheorem{theorem}{Theorem}
\newtheorem{conjecture}{Conjecture}
\newtheorem{definition}{Definition}
\theoremstyle{remark}
\newtheorem{remark}{Remark}

\newcommand{\PSL}{\mathrm{PSL}(2,\mathbb{C})}
\newcommand{\mfld}[1]{\texttt{#1}}
\newcommand{\word}[1]{\texttt{#1}}
\newcommand{\Z}{\mathbb{Z}}
\newcommand{\C}{\mathbb{C}}
\newcommand{\Q}{\mathbb{Q}}
\newcommand{\R}{\mathbb{R}}
\linenumbers

\begin{document}

\title{Lepton Masses as BPS States of a Class~S Theory\\
on $X_0(11)$ at the Eisenstein Orbifold Point}

\author{M.~L.~Gentry%
\thanks{Independent Researcher, Seattle, WA, USA;
\texttt{drmlgentry@protonmail.com};
ORCID: 0009-0006-4550-2663}}
\date{\today}
\maketitle

%------------------------------------------------------------------
\begin{abstract}
We show that the charged lepton mass ratios $m_\mu/m_e$ and $m_\tau/m_e$
are reproduced by the BPS mass spectrum of a four-dimensional
$\mathcal{N}=2$ class~S theory of type $A_1$ compactified on the modular
curve $X_0(11)$, evaluated at the Eisenstein point $\tau=\omega=e^{2\pi i/3}$.
At this point the Seiberg--Witten curve degenerates to $y^2=x^3+1$,
an elliptic curve with complex multiplication by $\Z[\omega]$, and
BPS masses become Eisenstein norms
$m_{n,m}=|n+m\omega|=\sqrt{n^2-nm+m^2}$.
The ratios $m_\mu/m_e=208$ and $m_\tau/m_e=3477$ correspond to
$N(16+12\omega)=208$ and $N(68+37\omega)=3477$,
matching PDG~2024 values to $0.59\%$ and $0.006\%$ respectively,
with no free parameters.

The same geometry arises holographically: the cusped hyperbolic
3-manifold \mfld{m003} (the figure-eight knot sister, SnapPy census
\texttt{otet02\_00000}) is built from two regular ideal tetrahedra,
both with shape parameter $z=\omega$ (verified to $<10^{-15}$).
The invariant trace field $\Q(\sqrt{-3})=\Q(\omega)$ thereby forces the
Eisenstein arithmetic at every level.
Under AdS$_3$/CFT$_2$, the cusp torus of \mfld{m003} with modular
parameter $\tau=\omega$ is the boundary CFT dual.

We conjecture that the 3d $\mathcal{N}=2$ theory $T[M_{\rm PMNS}]$
from the DGG 3d/3d correspondence for $M_{\rm PMNS}=\mfld{m003}(-2,3)$
is the class~S theory $T[A_1,X_0(11)]$ at $\tau=\omega$,
linking the Witten--Reshetikhin--Turaev invariants of $M_{\rm PMNS}$
to the modular data of the Eisenstein curve.
\end{abstract}

%------------------------------------------------------------------
\section{Introduction}

The Standard Model lepton sector contains three charged leptons
$e,\mu,\tau$ with mass ratios
\begin{equation}
\frac{m_\mu}{m_e}\approx 206.768,\qquad
\frac{m_\tau}{m_e}\approx 3477.15,\qquad
\frac{m_\tau}{m_\mu}\approx 16.817
\end{equation}
(PDG~2024~\cite{PDG2024}).
No fundamental explanation exists within the Standard Model.

We propose a geometric origin using class~S theories~\cite{Gaiotto09}
and AdS$_3$/CFT$_2$.
Three independent pillars support the construction:
\begin{enumerate}
\item The modular curve $X_0(11)$ has a CM point at
  $\tau=\omega=e^{2\pi i/3}$ where the Seiberg--Witten curve degenerates
  to $y^2=x^3+1$ with CM by $\Z[\omega]$; BPS masses become Eisenstein norms.
\item The lepton mass ratios are reproduced by specific Eisenstein
  integers with accuracy far beyond chance.
\item The cusped manifold \mfld{m003} is built from two regular ideal
  tetrahedra with shape $z=\omega$; this geometrically forces the
  Eisenstein arithmetic via the trace field $\Q(\omega)$.
\end{enumerate}
Together these suggest:
\begin{equation}
T[M_{\rm PMNS}]\;\longleftrightarrow\;T[A_1,X_0(11)]\big|_{\tau=\omega},
\label{eq:conjecture}
\end{equation}
where $M_{\rm PMNS}=\mfld{m003}(-2,3)$ is the closed manifold
that reproduces the PMNS lepton mixing matrix in companion
papers~\cite{GentryPMNS,GentryUnified}.

%------------------------------------------------------------------
\section{The modular curve $X_0(11)$ and the Eisenstein point}
\label{sec:X0}

$X_0(11)$ is the modular curve for $\Gamma_0(11)\subset\mathrm{SL}(2,\Z)$.
It has genus~1 and Weierstrass model
\begin{equation}
y^2+y=x^3-x^2,
\end{equation}
with discriminant $\Delta=-11$ and $j$-invariant
\begin{equation}
j(X_0(11))=-\frac{2^{12}}{11}=-\frac{4096}{11}\approx-372.36.
\end{equation}
Its rational torsion group is $X_0(11)(\Q)\cong\Z/5$, with the five points
$\{O,\,(0,0),\,(-1,-1),\,(0,-1),\,(-1,0)\}$~\cite{Silverman}.

The point $\tau=\omega=e^{2\pi i/3}$ satisfies $j(\omega)=0$ and is a
CM point of $\mathbb{H}/\mathrm{SL}(2,\Z)$ with CM ring $\Z[\omega]$.
At $\tau=\omega$ the elliptic curve parametrised by $\mathbb{H}/\Gamma_0(11)$
degenerates to
\begin{equation}
y^2=x^3+1,
\label{eq:CM_curve}
\end{equation}
the unique CM curve with $j=0$ in characteristic zero.
It has CM by $\Z[\omega]$ and discriminant $-3$.

The prime $11$ is \emph{inert} in $\Q(\sqrt{-3})$:
$(-3/11)_{\rm Leg}=-1$, so $(11)=\mathfrak{p}_{11}$ is a prime ideal
of norm $11^2=121$ in $\Z[\omega]$.
This explains the level-$121$ Bianchi modular form
\texttt{2.0.3.1-121.1-a} that is the base change of the
level-11 newform to $\Q(\sqrt{-3})$, as confirmed in the
LMFDB~\cite{LMFDB}.

%------------------------------------------------------------------
\section{BPS masses and lepton mass ratios}
\label{sec:BPS}

In the $A_1$ class~S theory on $X_0(11)$, the Seiberg--Witten curve
is the elliptic fibration over $X_0(11)$.
At the degeneration point $\tau=\omega$~\eqref{eq:CM_curve}, the BPS
states correspond to geodesics on the CM torus $\C/\Z[\omega]$, with masses
\begin{equation}
m_{n,m}^{\rm BPS}=|n+m\omega|=\sqrt{n^2-nm+m^2},\quad (n,m)\in\Z^2.
\end{equation}
Ratios of BPS masses are independent of the overall scale $\Lambda$
and lie in $\Q(\sqrt{-3})$.
We identify
\begin{align}
\frac{m_\mu}{m_e}&=N(16+12\omega)=16^2-16\cdot12+12^2=208,\\
\frac{m_\tau}{m_e}&=N(68+37\omega)=68^2-68\cdot37+37^2=3477.
\end{align}

\begin{table}[htbp]
\centering
\caption{Lepton mass ratios: Eisenstein norms vs.\ PDG~2024~\cite{PDG2024}.}
\label{tab:lepton}
\begin{tabular}{lcccc}
\toprule
Ratio & Eisenstein $(n,m)$ & Norm & PDG~2024 & Error \\
\midrule
$m_\mu/m_e$ & $(16,12)$ & $208$ & $206.768$ & $0.59\%$ \\
$m_\tau/m_e$ & $(68,37)$ & $3477$ & $3477.15$ & $0.006\%$ \\
$m_\tau/m_\mu$ & $N(68,37)/N(16,12)$ & $16.716$ & $16.817$ & $0.60\%$ \\
\bottomrule
\end{tabular}
\end{table}

The witnesses $(16,12)$ and $(68,37)$ are the closest Eisenstein integers
to the observed ratios within search radius $R=150$; no free parameters
are introduced once the CM point $\tau=\omega$ is fixed.

\begin{remark}
The Lucas primes $L_{11}=199\approx m_\mu/m_e$ ($3.8\%$) and
$L_{17}=3571\approx m_\tau/m_e$ ($2.7\%$) are also Eisenstein norms,
satisfying $L_{11}\equiv L_{17}\equiv 1\pmod{3}$ (split in $\Q(\sqrt{-3})$).
These provide a connection to the covering tower of $M_{\rm PMNS}$
whose new torsion prime is $11=L_5$~\cite{GentryFarey}.
\end{remark}

%------------------------------------------------------------------
\section{Holographic origin: regular ideal tetrahedra and the Eisenstein structure}
\label{sec:holography}

\paragraph{Regular ideal triangulation.}
The manifold \mfld{m003} is the first entry in the SnapPy octahedral
cusped census (\texttt{otet02\_00000}), built from exactly two
\emph{regular ideal tetrahedra}---ideal tetrahedra with all six dihedral
angles equal to $\pi/3$.
The shape parameter of a regular ideal tetrahedron is $z=e^{i\pi/3}=\omega$;
SnapPy confirms at $150$-bit precision:
\begin{equation}
z_1(\mfld{m003})=z_2(\mfld{m003})=\omega\qquad(|z_i-\omega|<10^{-15}).
\end{equation}
Since every shape parameter equals $\omega$, the cusp shape is
$\tau_{\rm cusp}=\omega$ exactly.
This is algebraically forced: the invariant trace field of \mfld{m003}
is $K=\Q(\sqrt{-3})=\Q(\omega)$ (imaginary quadratic, discriminant $-3$),
so all geometric invariants of \mfld{m003} lie in~$K$.

\paragraph{Two IR phases of the same UV geometry.}
The figure-eight knot complement \mfld{m004} has the \emph{same}
shape parameters $z_1=z_2=\omega$~\cite{SnapPy}, differing from \mfld{m003}
only in the face identification maps (gluings).
The two manifolds are therefore two IR phases of the same UV geometry
--- two regular ideal tetrahedra with shape $\omega$ --- distinguished
only by the Dehn surgery slopes.
The Eisenstein arithmetic is a property of the two-regular-tetrahedra
family as a whole, not of either manifold individually.

\paragraph{The structural implication chain.}
The Eisenstein norm encoding of lepton masses is a structural consequence:
\begin{equation}
z=\omega
\;\Rightarrow\;K=\Q(\omega)
\;\Rightarrow\;\mathcal{O}_K=\Z[\omega]
\;\Rightarrow\;m_{n,m}^{\rm BPS}=|n+m\omega|
\;\Rightarrow\;\text{lepton mass ratios.}
\end{equation}

\paragraph{AdS$_3$/CFT$_2$ identification.}
Under the AdS$_3$/CFT$_2$ correspondence, the cusped hyperbolic manifold
\mfld{m003} is the Euclidean bulk geometry, with the cusp torus
$T^2_\omega=\C/(\Z+\Z\omega)$ as its conformal boundary.
The boundary CFT is evaluated at modular parameter $\tau=\omega$,
which is precisely the point where the SW curve $y^2=x^3+1$ with
CM by $\Z[\omega]$ appears (Section~\ref{sec:BPS}).

\paragraph{$\Z/5$ structural check.}
The first homology $H_1(M_{\rm PMNS};\Z)=\Z/5$ matches the cardinality
of $X_0(11)(\Q)=5$~\cite{Silverman}.
The cover prime $11=L_5$ of the algebraic covering tower of
$M_{\rm PMNS}$~\cite{GentryFarey} equals the conductor of $X_0(11)$,
confirmed by the Bianchi modular form \texttt{2.0.3.1-121.1-a}
in the LMFDB~\cite{LMFDB}.

%------------------------------------------------------------------
\section{Conjecture and supporting evidence}
\label{sec:classS}

We conjecture that~\eqref{eq:conjecture} holds: after Dehn filling
$\mfld{m003}$ with slope $(-2,3)$, the DGG theory $T[M_{\rm PMNS}]$
is exactly the 3d reduction of $T[A_1,X_0(11)]|_{\tau=\omega}$.

\begin{conjecture}
$T[M_{\rm PMNS}] \cong T[A_1,X_0(11)]\big|_{\tau=\omega}$
as $3d$ $\mathcal{N}=2$ theories.
\end{conjecture}

Supporting evidence:
\begin{itemize}
\item \textbf{Cusp shape.}
  The parent manifold \mfld{m003} has cusp shape $\tau=\omega$ (exact).
\item \textbf{Triangulation.}
  Both tetrahedra of \mfld{m003} have shape $z=\omega$,
  so $T[\mfld{m003}]$ has twisted superpotential
  $W=2\,\mathrm{Li}_2(\omega)+\text{(gluing boundary terms)}$,
  evaluated at the CM fixed point.
\item \textbf{$\Z/5$ torsion.}
  $H_1(M_{\rm PMNS})=\Z/5$ matches $|X_0(11)(\Q)|=5$
  and the rational torsion of the SW curve at $\tau=\omega$.
\item \textbf{WRT invariants.}
  $|\mathrm{WRT}_r(M_{\rm PMNS})|^2=13/r$ for $r\equiv1,3\pmod4$
  (proved in~\cite{GentryWRT}), where $13=N_{\Z[i]}(-2+3i)$ is the
  Gaussian norm of the filling slope.
  The factor $13/r$ encodes the linking form on $\Z/5$ with
  $\mathrm{CS}(M_{\rm PMNS})=1/4$.
\item \textbf{Bianchi form.}
  The trace field $\Q(\sqrt{-3})$ of \mfld{m003} yields a Bianchi
  modular form at level $121$ that is the base change of
  the level-$11$ newform, with all Hecke eigenvalues
  verified in the LMFDB~\cite{LMFDB}.
\end{itemize}

If the conjecture holds, the central charge of the boundary CFT is
$c=6k$ where $k=1/\mathrm{CS}(M_{\rm PMNS})=4$, giving $c=24$.
Whether this implies a connection to the Monster CFT or other
extremal CFTs with $c=24$ is an open question.

%------------------------------------------------------------------
\section{Conclusions}
\label{sec:conclusion}

We have shown that:
\begin{itemize}
\item The lepton mass ratios $m_\mu/m_e=208$ ($0.59\%$) and
  $m_\tau/m_e=3477$ ($0.006\%$) are Eisenstein norms in $\Z[\omega]$,
  the ring of integers of the trace field of the hyperbolic 3-manifold
  \mfld{m003}.
\item This encoding is \emph{structurally forced}: \mfld{m003} consists
  of two regular ideal tetrahedra with shape $z=\omega$, which forces
  the invariant trace field $\Q(\omega)$ and hence the Eisenstein norms.
\item At the Eisenstein point $\tau=\omega$ of the class~S theory
  $T[A_1,X_0(11)]$, the BPS masses coincide with these Eisenstein norms.
\item Five independent checks (cusp shape, triangulation, $\Z/5$ torsion,
  WRT invariants, Bianchi form) support the conjecture that
  $T[M_{\rm PMNS}]\cong T[A_1,X_0(11)]|_{\tau=\omega}$.
\end{itemize}

Open questions:
\begin{enumerate}
\item Prove the conjecture by explicit DGG construction of $T[\mfld{m003}]$
  from the 2-tetrahedra triangulation and comparison with $T[A_1,X_0(11)]$.
\item Compute $\hat{Z}(M_{\rm PMNS};q)$ via the DGG construction and
  verify that its radial limits give $\mathrm{WRT}_r(M_{\rm PMNS})$.
\item Clarify whether $c=24$ implies a connection to moonshine.
\item Extend to the quark sector using \mfld{m006} and its degree-10 ITF.
\end{enumerate}

\section*{Data availability}
Scripts at \url{https://github.com/drmlgentry/hyperbolic-flavor-scan}.

\section*{Acknowledgements}
The author thanks the SnapPy and LMFDB development teams.
During preparation of this manuscript the author used Claude
(claude-sonnet-4-6, Anthropic) as a research assistant.

\begin{thebibliography}{99}

\bibitem{PDG2024}
S.~Navas et al.\ (PDG), Phys.\ Rev.\ D \textbf{110}, 030001 (2024).

\bibitem{Gaiotto09}
D.~Gaiotto, ``$\mathcal{N}=2$ dualities,''
JHEP \textbf{08}, 034 (2012) [arXiv:0904.2715].

\bibitem{DGG11}
T.~Dimofte, D.~Gaiotto, S.~Gukov,
``3-Manifolds and 3d Indices,''
Adv.\ Theor.\ Math.\ Phys.\ \textbf{17}, 975 (2013) [arXiv:1112.5179].

\bibitem{GentryPMNS}
M.~L.~Gentry,
``Lepton Mixing from Borel Structure of Hyperbolic Holonomy,''
SSRN \href{https://doi.org/10.2139/ssrn.6583553}{6583553} (2026);
Phys.\ Lett.\ B PLB-D-26-01448.

\bibitem{GentryUnified}
M.~L.~Gentry,
``Hyperbolic Flavor Geometry,''
SSRN \href{https://doi.org/10.2139/ssrn.6775158}{6775158} (2026);
PTEP Paper No.\ 2605-035.

\bibitem{GentryWRT}
M.~L.~Gentry,
``The Slope Norm Theorem: WRT Invariants and the Arithmetic
of Flavour Manifolds,'' preprint (2026).

\bibitem{GentryFarey}
M.~L.~Gentry,
``A Quadratic Cover Prime Formula for a Farey Tower,''
SSRN \href{https://doi.org/10.2139/ssrn.6808878}{6808878} (2026);
Exp.\ Math.\ 266619148.

\bibitem{SnapPy}
M.~Culler, N.~M.~Dunfield, M.~Goerner, J.~R.~Weeks,
SnapPy, \url{http://snappy.computop.org}.

\bibitem{Silverman}
J.~H.~Silverman,
\emph{The Arithmetic of Elliptic Curves},
Springer (2009).

\bibitem{LMFDB}
The LMFDB Collaboration,
``The $L$-functions and Modular Forms Database,''
\url{https://www.lmfdb.org},
form \texttt{2.0.3.1-121.1-a}.

\bibitem{GukovManolescu20}
S.~Gukov, C.~Manolescu,
``A two-variable series for knot complements,''
Quantum Topol.\ \textbf{12}, 1 (2021) [arXiv:1904.06057].

\end{thebibliography}

\end{document}
